% 기어식 기계의 앞모습 (개념도. 사진이나 도면이 아니다). 그림 fig:hw-gear-front
% 좌표는 mm, 원점은 바닥의 왼쪽 끝. 치수의 출처는 hardware/gear/layout.md 3.9--3.11절과
% hardware/gear/cad/params.py(BANK_W 1,880, DRIVE_W 360, PLINTH_H 400, ROW_Y 980/1,100/1,220,
% INDICATOR_D 70, 지시 드럼 축 = 구동부 시작 + 150). 대각선 보드의 잭 간격은 부록 E(13열 모두 110, 열 가운데 100 + 110 j).
\begin{tikzpicture}[x=0.07mm,y=0.07mm,
    lab/.style={font=\tiny,align=center},
    jack/.style={draw=corange,fill=corangebg,line width=0.3pt},
    tk/.style={cgray,line width=0.3pt}]
  % ---------- 바닥과 받침대 ----------
  \draw[cgray,line width=0.6pt] (-60,0) -- (1940,0);
  \draw[draw=cgray,fill=clight] (20,0) rectangle (1860,400);
  \node[lab,text=cgray] at (940,200) {받침대 (높이 400)};
  \draw[draw=cgray,fill=cgray!25] (0,400) rectangle (1880,450);
  % ---------- 뱅크 프레임 ----------
  \draw[draw=cgray,fill=clight!60,line width=0.8pt] (0,870) rectangle (1880,1350);
  % 열 자리 12개: 7--12번에 열 모듈, 1--6번은 막음판
  \foreach \s in {1,...,12} {
    \pgfmathsetmacro{\xl}{40+(\s-1)*120}
    \node[font=\tiny,text=cgray] at ({\xl+60},1322) {\s};
  }
  \foreach \s in {1,...,6} {
    \pgfmathsetmacro{\xl}{40+(\s-1)*120}
    \draw[draw=cgray,dashed,fill=white] ({\xl+2},900) rectangle ({\xl+118},1300);
  }
  \node[lab,text=cgray] at (400,1150) {1--6번: 막음판\\(나중에 채운다)};
  \foreach \s in {7,...,12} {
    \pgfmathsetmacro{\xl}{40+(\s-1)*120}
    \draw[draw=cgray,fill=white] ({\xl+2},900) rectangle ({\xl+118},1300);
    \foreach \y in {980,1100,1220} {
      \draw[draw=cgray,fill=cbluebg] ({\xl+60},\y) circle (48);
      \fill[cgray] ({\xl+60},\y) circle (6);
      \draw[cred,line width=0.8pt] ({\xl+60},{\y+48}) -- ({\xl+60},{\y+34});
    }
  }
  \node[lab,anchor=east] at (755,1220) {위 줄 (빠름)};
  \node[lab,anchor=east] at (755,1100) {가운데 줄};
  \node[lab,anchor=east] at (755,980) {아래 줄 (느림)};
  % ---------- 구동부 (오른쪽 끝, 폭 360) ----------
  \draw[draw=cgray,fill=white,line width=0.6pt] (1480,900) rectangle (1840,1300);
  \draw[draw=cgray,fill=clight,pattern=north east lines,pattern color=cgray!50] (1680,920) rectangle (1830,1280);
  \node[lab,fill=white,inner sep=1pt] at (1755,1100) {기어 창\\(막힘)};
  \draw[draw=cgray,fill=cyellowbg!60] (1595,940) rectangle (1665,1260);
  \foreach \y in {980,1100,1220} {
    \draw[draw=cgray,fill=cyellowbg] (1630,\y) circle (35);
    \node[font=\tiny\ttfamily] at (1630,\y) {K};
  }
  \node[lab,anchor=east] at (1590,1260) {지시\\드럼};
  \node[lab,text=cgray] at (1660,1318) {구동부};
  % ---------- 화면 (구동부 위) ----------
  \draw[draw=cgray,fill=white] (1555,1355) -- (1765,1355) -- (1745,1540) -- (1575,1540) -- cycle;
  \node[lab] at (1660,1450) {화면\\(10인치급)};
  % ---------- 투명 덮개 ----------
  \draw[cblue,dashed,line width=0.7pt,rounded corners=2pt] (30,885) rectangle (1850,1340);
  \node[lab,text=cblue,anchor=south west] at (40,1350) {투명 앞 덮개 (위쪽 경첩, 열면 모터 전원이 끊긴다)};
  % ---------- 패널 (폭 1,800, 높이 420) ----------
  \draw[draw=cgray,fill=white,rounded corners=2pt,line width=0.6pt] (40,450) rectangle (1840,870);
  \draw[cgray,dashed,line width=0.4pt] (1480,455) -- (1480,865);
  \node[lab,text=cgray,anchor=north west] at (45,866) {메뉴 배선 패널};
  \node[lab,text=cgray,anchor=north east] at (1835,866) {조작부};
  % 유닛 잭: 열마다 IN(위)과 OUT(아래)
  \foreach \s in {1,...,12} {
    \pgfmathsetmacro{\xc}{40+(\s-0.5)*120}
    \draw[jack] ({\xc-35},812) rectangle ({\xc+35},840);
    \draw[jack] ({\xc-35},778) rectangle ({\xc+35},806);
  }
  % 대각선 보드 다섯 장과 글자 잭(열마다 두 글자 × 잭 둘)
  \foreach \xa/\xb/\n in {57.75/389.33/D1, 391.75/719.34/D2, 721.75/1049.34/D3, 1051.76/1269.34/D4, 1271.76/1462.25/D5} {
    \draw[draw=cgray,fill=clight!70] (\xa,640) rectangle (\xb,770);
    \node[font=\tiny,text=cgray] at ({(\xa+\xb)/2},628) {\n};
  }
  \foreach \x in {100,210,320,430,540,650,760,870,980,1090,1200,1310,1420} {
    \foreach \y in {655,685,715,745} { \draw[jack] ({\x-18},{\y-9}) rectangle ({\x+18},{\y+9}); }
  }
  % 표시 릴레이 26개(투명 케이스)와 레지스터 램프 26개
  \draw[draw=cgray,fill=white] (300,585) rectangle (1080,620);
  \foreach \i in {0,...,25} { \draw[draw=cgray,fill=cyellowbg] ({305+\i*30},589) rectangle ++(20,27); }
  \foreach \i in {0,...,25} { \draw[draw=cgray,fill=cyellowbg] ({100+\i*55},530) circle (11); }
  \node[font=\tiny\ttfamily,text=cgray] at (100,500) {A};
  \node[font=\tiny\ttfamily,text=cgray] at (1475,500) {Z};
  \node[lab,anchor=west] at (1090,602) {표시 릴레이 26};
  \node[lab] at (790,500) {레지스터 램프 26};
  % 조작부: 비상 정지(바닥에서 약 800), Start·Stop·Home, 카드 판독기, 분기 잭
  \draw[draw=cred,fill=credbg,line width=0.7pt] (1540,800) circle (24);
  \node[font=\tiny,anchor=north] at (1540,774) {비상 정지};
  \foreach \x/\t/\c in {1625/시작/cgreen, 1700/정지/cgray, 1775/Home/cgray} {
    \draw[draw=\c,fill=white] ({\x-30},786) rectangle ({\x+30},814);
    \node[font=\tiny] at (\x,800) {\t};
  }
  \draw[draw=cgray,fill=white] (1620,690) rectangle (1830,738);
  \node[lab] at (1725,714) {카드 판독기};
  \draw[draw=cgray,fill=clight] (1490,465) rectangle (1830,625);
  \foreach \r in {0,...,3} { \foreach \c in {0,...,3} {
    \draw[jack] ({1520+\c*85-18},{490+\r*33-9}) rectangle ++(36,18); } }
  \node[lab,anchor=south] at (1660,625) {분기 잭 4조 (패치 패널)};
  % ---------- 높이의 눈금 (바닥에서, mm) ----------
  \foreach \h/\t in {0/0, 400/400, 450/450, 870/870, 980/980, 1100/{1,100}, 1220/{1,220}, 1350/{1,350}} {
    \draw[tk] (-45,\h) -- (-5,\h);
    \node[font=\tiny,text=cgray,anchor=east] at (-48,\h) {\t};
  }
  % ---------- 폭의 치수 ----------
  \foreach \x in {0,40,1480,1840,1880} { \draw[tk] (\x,-15) -- (\x,-175); }
  \draw[<->,cgray,line width=0.4pt] (40,-60) -- (1480,-60) node[midway,fill=white,font=\tiny] {1,440 = 12 자리 × 120};
  \draw[<->,cgray,line width=0.4pt] (1480,-60) -- (1840,-60) node[midway,fill=white,font=\tiny] {구동부 360};
  \draw[<->,cgray,line width=0.4pt] (0,-140) -- (1880,-140) node[midway,fill=white,font=\tiny] {1,880 = 40 + 1,440 + 360 + 40};
\end{tikzpicture}
